\subsection{Extensibility offered by \enumo}
\label{sec:case-study-llm}
We have demonstrated that \enumo's core operators enable more
  flexible and incremental rule inference strategies than prior work.
In this section, we show how to customize and extend the \enumo DSL to
  support entirely new techniques for rule synthesis.
Given the increasing popularity of large language models (LLMs) for synthesis
  tasks~\citep{hysynth,asap,shraddha,paper1,paper2,paper3,paper4},
  we extend \enumo to support rule inference using LLMs.
For \eqsat applications, rule soundness is paramount: even a single
  unsound rule can propagate merges through the \egraph.
As such, LLM-generated rulesets must be treated as untrusted and
  should not be used directly without processing to eliminate unsound rules.
Since \enumo already includes operators for processing untrusted rule candidates
  (\C{minimize} and \C{is_valid}), we find it to be a natural fit for
  working with untrusted LLM-generated rulesets.
Concretely, we added two new operators to \enumo{}: \C{workload_from_llm} and
  \C{ruleset_from_llm}.
The input to each operator is a natural-language prompt string\footnote{
  The prompts sent to the models are under \C{jfp/} at tag
  \C{jfp2026} of the repository
  (\url{https://github.com/uwplse/ruler/tree/jfp2026}).
}
  and the output is a \C{Workload} or \C{Ruleset}, respectively.
We measure how easily LLMs can be incorporated into \enumo and how the
  resulting operators compose with the existing ones.
Comparing the results from different LLMs is beyond the scope of this work,
  as is comparing LLM performance with different prompts for the same task.
To reduce variance of LLM performance, we use several models\footnote{
  Gemini 3.6 Flash, GPT 5.6 Luna, and Claude Sonnet 5}, querying
  each model twice, and aggregating all of the responses together.
These two extensions to \enumo required only 249 lines of Rust
  code to implement\footnote{The LLM client module plus the two
  operators' constructors and the workload's term filter, at the
  tag cited above.}, including only minimal changes to existing
  code.

\subsubsection{Ruleset synthesis}
In the first case study, we use LLMs for end-to-end ruleset synthesis
  in three domains: Halide, Trigonometric, and Exponential.
In each case, our prompt
  explained the task of rule inference at a high level
  and specified the valid values and operators in the domain
  (the domain descriptions in \autoref{table:llmrulesets}).
The prompt also included example rules and listed classes of rules to cover:
  for example, commutativity, associativity, and distributivity for Halide;
  angle-sum and double-angle identities for Trigonometric;
  and product, quotient, and power laws for Exponential.
Finally, it stressed that every rule must be sound wherever both sides
  are defined and asked for plain text, one rule per line, with no commentary.
We parse rules out of the LLM response, discarding any malformed
  rules, and treat the resulting ruleset as \textit{untrusted rule candidates},
  which we verify as explained below.
Since we aggregate responses from several models and queries,
  there are many redundant rules generated, which
  we eliminate using \enumo's \C{minimize} operator.

\paragraph*{Verifying LLM-generated rule candidates.}
For Halide, we validate the LLM-generated rule candidates using
  \enumo's \C{is_valid} operator,
  which checks the rule by encoding it as an SMT query.
Empirically, we find that unsound rules from LLMs are a genuine problem:
  in the Halide domain, 36 of the 977 LLM-generated rule candidates were
  found to be unsound.
While this is a relatively small percentage of the rules, in \eqsat settings,
  the presence of a single unsound rule can invalidate all of the
  equalities represented in the \egraph.
For Trigonometric and Exponential, validity
  is undecidable in general~\citep{dreal},
  so SMT-based validation does not apply.
Instead, we reuse an existing \enumo operator, \C{can_derive},
  to check soundness by derivability,
  following the same intuition as our \ff algorithm (\autoref{sec:lift}).
Specifically, we attempt to derive each
  LLM-synthesized rule candidate using
  the known rules for each domain: \C{R}\footnote{
    While unsound rules were acceptable in \autoref{sec:lift} where we could ensure
    that no undefined terms were represented in the \egraph, in this setting,
    we exclude the 13 arithmetic rules involving division that are unsound at zero.
  } \C{+ C + E} for Trigonometric and \C{RAT + EXP}
  for Exponential (\autoref{sec:lift}).
Assuming soundness of these trusted rules,
  if a rule candidate can be derived
  from them,
  then the rule candidate is sound.
If a rule candidate cannot be derived, it is filtered out,
  even though it may be sound.

\paragraph*{Eliminating redundant LLM-generated rules.}
Once we have filtered out potentially unsound rule candidates,
  we address the problem
  of redundant rules.
We use \enumo's \C{minimize} operator (\autoref{subsec:ruleset-sem}),
  which reduces a ruleset with respect to a set of prior rules.
For Halide, we consider three sets of prior rules: \C{None} (the empty ruleset),
  \C{A5} (rules synthesized by exhaustively enumerating terms up to size 5),
  and \C{ENUMO} (rules synthesized using a custom \enumo
  program, \autoref{subsec:halideeval}).
For Trigonometric, we use the arithmetic rules \C{R} without the
  division rules (as in \autoref{table:herbie}) as prior rules for
  minimization; for Exponential, the rational rules \C{RAT} alone.

\begin{table}[t]
\resizebox{\textwidth}{!}{%
\begin{tabular}{|l|l|l|l|c|l|}
\hline
Ruleset Name & Domain                         & Domain Description                                                                                                                                                                                 & Prompt Rules & \multicolumn{1}{l|}{Validation}                                                     & Prior Rules         \\ \hline
LLM-1        & \multirow{6}{*}{Halide}        & \multirow{6}{*}{\begin{tabular}[c]{@{}l@{}}Values: integers\\ Unary Operators: -, !\\ Binary Operators: \C{<}, \C{<=}, ==, !=, \&\&, \C{||}, \^{}, +, -, *, min, max\\ Ternary Operators: select\end{tabular}} & -            & \multirow{6}{*}{SMT}                                                                & -                   \\
LLM-2        &                                &                                                                                                                                                                                               & LLM-1        &                                                                                     & LLM-1               \\
LLM-A5-1     &                                &                                                                                                                                                                                               & -            &                                                                                     & A5                  \\
LLM-A5-2     &                                &                                                                                                                                                                                               & LLM-A5-1     &                                                                                     & A5 + LLM-A5-1       \\
LLM-ENUMO-1  &                                &                                                                                                                                                                                               & -            &                                                                                     & ENUMO               \\
LLM-ENUMO-2  &                                &                                                                                                                                                                                               & LLM-ENUMO-1  &                                                                                     & ENUMO + LLM-ENUMO-1 \\ \hline
LLM-1        & \multirow{3}{*}{Trigonometric} & \multirow{3}{*}{\begin{tabular}[c]{@{}l@{}}Values: real numbers, PI\\ Unary operators: -, sin, cos, tan, sqr\\ Binary operators: +, -, *, /\\\end{tabular}}                                   & -            & \multirow{3}{*}{\begin{tabular}[c]{@{}c@{}}Derivability from\\ R + C + E\end{tabular}} & R                  \\
LLM-2        &                                &                                                                                                                                                                                               & LLM-1        &                                                                                     & R + LLM-1          \\
             &                                &                                                                                                                                                                                               &              &                                                                                     &                     \\ \hline
LLM-1        & \multirow{3}{*}{Exponential}   & \multirow{3}{*}{\begin{tabular}[c]{@{}l@{}}Values: real numbers\\ Unary operators: -, exp, log, sqrt, cbrt\\ Binary operators: +, -, *, /, pow\\\end{tabular}}                                  & -            & \multirow{3}{*}{\begin{tabular}[c]{@{}c@{}}Derivability from\\ RAT + EXP\end{tabular}} & RAT                  \\
LLM-2        &                                &                                                                                                                                                                                               & LLM-1        &                                                                                     & RAT + LLM-1          \\
             &                                &                                                                                                                                                                                               &              &                                                                                     &                     \\ \hline
\end{tabular}%
}
\caption{Description of how each ruleset is generated using LLMs.}
\label{table:llmrulesets}
\end{table}

In this experiment, we use LLMs
  in an iterative manner to
  incrementally construct rulesets.
After generating rules, we prompt
  the LLMs again in order to find missing rules.
The second prompt restates the domain, lists the rules already in the
  ruleset (the prompt rules in \autoref{table:llmrulesets}),
  and asks for sound rules missing from that list, excluding trivial
  variants such as renamed variables or swapped arguments of
  commutative operators.
This results in a new set of rule candidates,
  which we validate and minimize again
  for each domain.
For this second round of minimization,
  the rules generated in the first round
  are also included as prior rules.
\autoref{table:llmrulesets} summarizes the rulesets generated for this case study.

For each ruleset, we measured the
  \lhsandrhs derivability compared to two baselines: \C{ENUMO},
  the rules synthesized using \enumo programs as
  outlined in \autoref{subsec:halideeval} and \autoref{subsec:eval-ff},
  and \C{HALIDE} (for Halide) or \C{HERBIE} (for Exponential and Trigonometric).
As in \autoref{table:herbie}, every deriving ruleset in the Trigonometric
  and Exponential matrices is extended with \C{R} or \C{RAT}, respectively;
  the LLM rulesets were minimized against these rules and are not
  meaningful without them.
For the Halide domain, we also compare to \C{A5} as a third baseline.
The full results are shown in \autoref{table:llm1}.

As expected, we find that LLMs are powerful tools for ruleset synthesis.
Measured against the expert-written \C{HALIDE} and \C{HERBIE} baselines,
  the best LLM-synthesized ruleset outperforms guided search in
  the Halide and Exponential domains.
However, in the Trigonometric domain, the LLM-synthesized ruleset underperforms
  compared to \C{ENUMO}.
In all domains, neither technique subsumes the other:
  \C{ENUMO} rulesets contain
  rules not derivable from the LLM rulesets, and vice versa,
  indicating that they find substantively different rules.

We also find that LLMs compose well with traditional rule inference techniques.
For example, \C{LLM-A5-1} outperforms both \C{A5} and \C{LLM-1}
  against the \C{ENUMO} baseline, suggesting that \textit{composing
  traditional rule inference methods with LLMs can lead to
  better rulesets than using either technique alone}.
Further, re-prompting shows that the new operators compose with
  themselves and with the rest of the pipeline: the minimized first-round
  ruleset is embedded in the second prompt, and the second round's
  candidates are validated and minimized against the first round's rules.

\begin{table}[ht]
    \centering

    \resizebox{\textwidth}{!}{%
    \begin{tabular}{|lr|lllllllll}
    \hline
    Ruleset      & Synthesis Time (s)  & HALIDE       & A5           & ENUMO                             & LLM-1        & LLM-2       & LLM-A5-1    & LLM-A5-2    & LLM-ENUMO-1  & \multicolumn{1}{l|}{LLM-ENUMO-2} \\ \hline
    A5           & 15.7\phantom{$^{+}$} & 45.1\% (10.6) & -             & 69.3\% (26.5)                     & 53.0\% (8.1)  & 53.8\% (8.0) & 4.4\% (4.5)  & 16.5\% (5.5) & 14.9\% (3.5) & \multicolumn{1}{l|}{23.8\% (4.4)} \\
    ENUMO        & 49.2\phantom{$^{+}$} & 90.9\% (22.1) & 88.3\% (56.6) & -                                  & 74.5\% (32.6) & 74.9\% (34.6) & 44.4\% (15.5) & 48.5\% (18.7) & 17.9\% (21.2) & \multicolumn{1}{l|}{31.2\% (18.7)} \\ \cline{6-11}
    LLM-1        & 1043.7\phantom{$^{+}$} & 93.5\% (5.6)  & 64.0\% (19.8) & \multicolumn{1}{l|}{65.0\% (38.7)} &               &              &              &              &               &                                  \\
    LLM-2        & 326.8$^{+}$         & 94.3\% (4.7)  & 67.7\% (19.9) & \multicolumn{1}{l|}{68.7\% (39.5)} &               &              &              &              &               &                                  \\
    LLM-A5-1     & 1043.5\phantom{$^{+}$} & 90.1\% (38.1) & -             & \multicolumn{1}{l|}{81.2\% (25.3)} &               &              &              &              &               &                                  \\
    LLM-A5-2     & 586.2$^{+}$         & 89.4\% (36.5) & -             & \multicolumn{1}{l|}{81.3\% (17.2)} &               &              &              &              &               &                                  \\
    LLM-ENUMO-1  & 1043.5\phantom{$^{+}$} & 91.2\% (24.6) & 91.5\% (11.8) & \multicolumn{1}{l|}{-}             &               &              &              &              &               &                                  \\
    LLM-ENUMO-2  & 734.7$^{+}$         & 94.2\% (22.9) & 90.8\% (12.9) & \multicolumn{1}{l|}{-}             &               &              &              &              &               &                                  \\ \cline{1-5}
    \end{tabular}%
    }

    \vspace{1em}

    \begin{minipage}{0.45\textwidth}
        \centering
        \resizebox{\textwidth}{!}{%
        \begin{tabular}{|lr|llll}
        \hline
        Ruleset     & Synthesis Time (s)  & HERBIE       & ENUMO                             & LLM-1        & \multicolumn{1}{l|}{LLM-2}        \\ \hline
        ENUMO       & 4.6\phantom{$^{+}$} & 43.9\% (2.7) & -                                 & 89.3\% (1.0) & \multicolumn{1}{l|}{90.6\% (1.0)} \\ \cline{5-6}
        LLM-1       & 516.0\phantom{$^{+}$} & 52.4\% (1.7) & \multicolumn{1}{l|}{65.0\% (0.5)} &              &                                   \\
        LLM-2       & 399.3$^{+}$         & 53.7\% (1.5) & \multicolumn{1}{l|}{70.0\% (0.5)}  &              &                                   \\ \cline{1-4}
        \end{tabular}%
        }
    \end{minipage}
    \hfill
    \begin{minipage}{0.45\textwidth}
        \centering
        \resizebox{\textwidth}{!}{%
        \begin{tabular}{|lr|llll}
        \hline
        Ruleset  & Synthesis Time (s)  & HERBIE        & ENUMO                             & LLM-1        & \multicolumn{1}{l|}{LLM-2}        \\ \hline
        ENUMO    & 692.4\phantom{$^{+}$} & 35.6\% (14.0) & -                                 & 72.7\% (1.4) & \multicolumn{1}{l|}{75.0\% (1.4)} \\ \cline{5-6}
        LLM-1    & 464.3\phantom{$^{+}$} & 13.3\% (1.5)  & \multicolumn{1}{l|}{44.6\% (1.1)} &              &                                   \\
        LLM-2    & 770.0$^{+}$         & 13.3\% (1.7)  & \multicolumn{1}{l|}{44.6\% (1.1)}  &              &                                   \\ \cline{1-4}
        \end{tabular}%
        }
    \end{minipage}

    \caption{Derivability matrices for our LLM rule synthesis case study.
    Halide (Top), Exponential (Left), and Trigonometric (Right).
    Derivability times in seconds are shown in parentheses.
    Reprompted rulesets show only the second round synthesis times
    (marked with $^{+}$) and do not include the first
    round's synthesis time.
    Cell values show the \lhsandrhs derivability results of using the
    row ruleset to derive the column ruleset.
    For Exponential and Trigonometric, the row ruleset is extended with
    \C{RAT} or \C{R} before deriving, as in \autoref{table:herbie}.
    A dash marks cells where the column ruleset is contained in the row ruleset,
    so derivability holds by construction.
    Derivability between LLM rulesets is not included because
    the goal is not to compare LLM-synthesized rulesets against each other.
    }
    \label{table:llm1}
\end{table}

\subsubsection{Workload synthesis}
\label{subsec:case-study-workload}
The second case study explores the use of LLMs for term enumeration.
Rather than prompting the model for complete rulesets, we prompt the model
  to enumerate terms from the domain to make an \enumo workload.
For this case study, we focus on the Halide domain.
The prompt specifies the constants, variables, and operators that terms may use
  and gives several example terms.
It also explains that
  rules will be inferred from pairs of equivalent terms in the workload,
  so the terms should come in clusters that are likely to be equivalent,
  vary in size and nesting depth, and should cover all of the operators and their combinations.
The \enumo \C{workload_from_llm} operator then constructs an \enumo \C{Workload}
  from the valid terms in the LLM response,
  skipping any malformed s-expressions
  or invalid terms.
We also filter out any terms that use variables other than those specified
  in order to keep the cvec size manageable\footnote{cvec length grows exponentially
   with the number of variables in the \C{Workload}.}.

For the Halide domain,
  our prompt tells each model to generate at least 1000 terms.
Across two queries per model, the responses contained
  3407 (Gemini 3.6 Flash), 3195 (GPT 5.6 Luna), and 2111 (Claude Sonnet 5) terms,
  which aggregated to 5952 unique terms, 5878 of them valid.
In total, generating this workload took 973.3 seconds.

To synthesize rules from the workload,
  we follow the typical \enumo pipeline:
  convert the workload to an \egraph,
  compress the \egraph using prior rules,
  find rule candidates using cvec matching, and
  minimize the rule candidates
  to eliminate redundant and unsound rules.
The LLM-synthesized workload composes seamlessly into the
rest of the rule inference pipeline, enabling entirely new
strategies for rule inference without modifying any of \enumo's other operators.

We consider four sets of prior rules:
  \C{None}, \C{A5}, \C{ENUMO} as in
  the first case study, and also
  \C{LLM-2}, which is the LLM-generated ruleset from the
  first case study for Halide.
For each ruleset, we measured the \lhsandrhs derivability compared to
  three baselines: \C{HALIDE}, \C{A5}, and \C{ENUMO}.
The full results are shown in \autoref{table:llm2}.

As in the first case study, we find that incorporating LLMs can lead to more
  powerful rulesets than using either technique alone.
Inferring rules directly from the LLM-generated workload yields 2181 rules,
  but is quite slow (roughly 1.5 hours), with the majority of the time
  spent minimizing nearly 200,000 rule candidates.
Furthermore, the resulting ruleset is quite weak, deriving only 28.1\% of
  the \C{HALIDE} baseline.
However, it is fast and easy to use \enumo to generate a set of rules
  by exhaustively enumerating terms up to size 5 (\C{A5}).
If we use these rules as a starting point for rule inference using \C{LLM-W},
  we drastically reduce the number of candidates under consideration to only 7195,
  and minimization finishes in under a minute, resulting in 369 rules.
Perhaps surprisingly, these rules are also much more powerful than
  the directly-synthesized ruleset, deriving 69.7\% of \C{HALIDE}'s.
This suggests a promising strategy for workload-driven rule inference:
  use traditional techniques for small rules, where exhaustive and guided search
  are fast and reliable, then use LLMs to sample terms that are harder to find using search.

Notably, \C{LLM-W-LLM-2}, which was entirely synthesized by LLMs,
  with no domain expert guidance,
  outperforms \C{ENUMO} against the \C{HALIDE} baseline,
  suggesting that incorporating LLMs into rule inference pipelines could
  lower the barrier to producing high-quality rulesets.

\begin{table}[h]
\footnotesize
\resizebox{\textwidth}{!}{%
\begin{tabular}{|l|llllll|}
\hline
Ruleset     & \multicolumn{1}{l}{\# Prior Rules} & \multicolumn{1}{l}{\# New Rules} & \multicolumn{1}{r}{Rule Synthesis Time (s)} & HALIDE      & A5          & ENUMO        \\ \hline
A5          & 0             & 480                         & 15.7\phantom{$^{+}$}         & 45.1\% (10.6) & -             & 69.3\% (26.5) \\
ENUMO       & 0             & 840                         & 49.2\phantom{$^{+}$}         & 90.9\% (22.1) & 88.3\% (56.6) & -              \\
LLM-2       & 0             & 370                         & 1370.4\phantom{$^{+}$}       & 94.3\% (4.7)  & 67.7\% (19.9) & 68.7\% (39.5)  \\
LLM-W       & 0             & 2181                        & 5510.8$^{+}$                 & 28.1\% (3291.6)  & 35.8\% (2359.6)  & 25.4\% (6965.6)   \\
LLM-W-A5    & 480           & 369                         & 49.1$^{+}$                   & 69.7\% (1006.6) & -             & 74.4\% (214.8)  \\
LLM-W-ENUMO & 840           & 332                         & 74.9$^{+}$                   & 91.6\% (783.1) & 90.4\% (608.7) & -              \\
LLM-W-LLM-2 & 370           & 200                         & 27.9$^{+}$                   & 93.1\% (40.9)  & 78.1\% (48.4) & 73.9\% (99.8) \\ \hline
\end{tabular}%
}
\caption{
  Derivability comparison between rulesets synthesized from an LLM-generated workload.
  Cell values show the \lhsandrhs derivability results of using the
    row ruleset to derive the column ruleset.
  Each row ruleset consists of its prior rules together with the new rules
    synthesized on top of them, and derivability is measured using both.
  Derivability times in seconds are shown in parentheses.
  Dashes indicate cells where the column ruleset is contained within the row ruleset,
    and derivability is guaranteed by construction.
  The time to synthesize \texttt{LLM-2} includes the LLM query time.
  The synthesis times for the \texttt{LLM-W} rulesets (marked with $^{+}$) do not include the time to
  generate the workload from the LLMs, which was 973.3s.}%
\label{table:llm2}
\end{table}

 These two case studies show the extensibility of
 \enumo as new rule inference techniques emerge.
Adding support for LLMs to the \enumo implementation was straightforward,
  requiring minimal changes to existing code.
The new operators compose correctly with existing \enumo operators,
  as shown in \autoref{subsec:case-study-workload}, where we use the
  new \C{workload_from_llm} operator to construct a workload, leaving
  the rest of the rule inference pipeline unchanged.
The operators can also be composed with existing operators in novel ways,
  as in the Trigonometric and Exponential case studies, where we
  use \enumo's \C{can_derive} operator to determine validity of the LLM-generated
  rule candidates.

\paragraph*{Threats to validity.}
Our claims concern the ease of extending \enumo and the composability of
  the resulting operators, not the performance of LLMs relative to guided
  search, which varied widely across domains (\autoref{table:llm1}).
In these case studies, we did not attempt to
  compare performance between LLM models,
  nor did we rigorously explore how prompt
  engineering impacts the resulting rulesets.
Both are interesting directions of future work.
Due to the nondeterministic nature of LLMs,
  the specific rules and terms generated differ from run to run.

\subsection{Ruleset manipulation with \enumo}
\label{sec:case-study-bv}

\begin{table}[h]
\footnotesize
\begin{tabular}{llll}
Domain  & Generated Rules (Time)      & Valid BV4 Rules (Time)     & Validated $\rightarrow$ Generated \\ \cline{1-4}
BV8     & 230 (32.78)                 &  230 (3.19)                & (100\%, 100\%) \\
BV16    & 236 (85.50)                 &  224 (8.36)                & (97\%, 97\%) \\
BV32    & 232 (191.74)                &  224 (10.81)               & (98\%, 99\%) \\
BV64    & 250 (43.70)                 &  220 (16.97)               & (93\%, 93\%) \\
BV128   & 190 (1784.14)               &  210 (38.68)               & (90\%, 91\%) \\
\end{tabular}%
\caption{
  Comparison of rule synthesis for different widths of bitvectors.
  Shown for each bitvector width are
    (i) the number of rules generated from an
    \enumo program (time in seconds) for that domain,
    (ii) the number of \enumo-synthesized BV4 rules
    that are valid in that domain (time in seconds), and
    (iii) the percentage of the generated rules
    that are derivable from the validated BV4 rules
    ( \lhs and \lhsandrhs derivability, in that order).
}%
\label{table:bv}
\end{table}

In this section, we show a case study in leveraging \enumo's operators for
  ruleset manipulation.
Here, we are interested in synthesizing rewrite rules over 4-bit bitvectors
  (BV4), and transforming those rules into a usable ruleset over
  larger bitvectors.
Synthesizing rules for small bitvectors is very fast because there are
  relatively few possible values in the domain.
Rules that work for small bitvectors are likely, but not guaranteed, to be
  valid for large bitvectors as well.
In this case study, we start by synthesizing BV4 rules using
  the same \enumo program as described in \autoref{subsec:eval-enumo}.
Then we use \enumo operators to translate the rules into the domain of larger
  bitvectors (BV8, BV16, BV32, BV64, and BV128).
Finally, we validate the rules in the new domain to find the subset of sound BV4
  rules that are still sound for larger bitvectors.
We compare these rules against rules that were synthesized directly in the
  larger bitvector domains using the same \enumo program as we used to
  synthesize BV4 rules.
The results are shown in \autoref{table:bv}.
Validating BV4 rules is much
  faster than synthesizing rules from scratch
  (38 seconds vs.\ 29 minutes for BV128)
  and still yields good results.
Across all bitvector sizes, the validated BV4 rules retained at least 90\% of
  the proving power of the directly generated rules, and in the case of
  8-bit bitvectors, the validated BV4 rules had equal proving power.

